% ============================= Introduction =============================
\section{Introduction}
\label{sec:intro}

Vision-language models now answer questions about \emph{where} objects are with
remarkable skill: on VSR~\cite{liu2023vsr}, a frozen 7B model answers $80.9\%$ of
test statements correctly, and on SITE~\cite{wang2025site} the official
spatial-relationship category reaches $59.2\%$ chance-adjusted accuracy, on par
with published open-source baselines. But there is a subfamily of spatial
relations on which every model we test stagnates: \emph{object-intrinsic
orientation}. ``The person is facing away from the horse,'' ``the chair is parallel
to the table,'' ``the cat is facing the dog'' --- these statements probe whether a
model can read a direction off an object's own geometry, not merely compare two
positions. We find that this capability is remarkably resistant to improvement:
across two model families, five fine-tuning conditions, four representation
probes, an explicit two-stage reasoning pipeline, a logical-consistency analysis,
and an external benchmark, the orientation bottleneck persists.

The contribution of this paper is not another method; it is a
\emph{decomposition} of the failure that rules in and out candidate mechanisms,
all within one evaluation protocol. Concretely:

\begin{itemize}
\item \textbf{Scaling does not fix it.} Orientation is the weakest family of the
eight VSR relation families for both models ($62.8\%$ for 2B, $63.5\%$ for 7B vs.\
$75$--$93\%$ for easy families), and 2B$\rightarrow$7B scaling moves it by less than
a point.
\item \textbf{The representation itself is a bottleneck.} Linear, MLP, patch-vote,
and object-grounded probes trained on frozen Qwen2-VL-7B vision features extract
facing/facing-away at or near the majority baseline, even when readouts are
conditioned on the subject and reference object regions.
\item \textbf{Vision-side adaptation and explicit decomposition fail.}
Adapting the patch-merger or the upper vision blocks with LoRA does not change
orientation accuracy ($64.2\%$ vs.\ $65.7\%$, both conditions \emph{significantly
worse} overall, $p \leq 0.012$), and a ``localize-then-classify'' pipeline using
grounded boxes plus geometry features is significantly \emph{worse} than the
generative control ($55.1$--$58.3\%$ vs.\ $65.7\%$, all $p \leq 0.004$).
\item \textbf{Coherence and accuracy are separable.} The zero-shot 7B model
contradicts itself on $63.1\%$ of facing/facing-away statement pairs. LM-side
training repairs this, and hard-negative training raises facing-pair consistency
from $36.9\%$ to $77.7\%$ while accuracy is \emph{identical} ($68.9\%$).
\item \textbf{The weakness generalizes, and is specific.} On a preregistered
external validation over 2{,}591 SITE images, orientation-vocabulary questions
collapse to $48.3\%$ raw / $24.0\%$ chance-adjusted accuracy while the official
spatial-relationship category stays at $75.1\%$ / $59.2\%$. VSR images embedded in
SITE score $85.9\%$, ruling out a trivial image-quality explanation.
\end{itemize}

These results give the community a precise, well-measured target ---
\emph{object-intrinsic direction} --- that no single intervention we tried can
move, plus a template for decomposing any claimed spatial failure into its
representational and behavioral components. All numbers in this paper are traced
to committed artifacts (App.~G).
